\documentclass[10pt,a4paper]{article}
\usepackage[T1]{fontenc}
\usepackage[utf8]{inputenc}
\usepackage{lmodern}
\usepackage{amsmath,amssymb}
\usepackage[margin=21mm]{geometry}
\usepackage{microtype}
\usepackage[hidelinks]{hyperref}
\usepackage{enumitem}
\setlist[itemize]{leftmargin=*,nosep,topsep=3pt}
\setlength{\parindent}{0pt}
\setlength{\parskip}{5pt}
\setlength{\emergencystretch}{3em}
\newcommand{\question}[1]{\par\medskip{\normalsize\bfseries #1}\par\smallskip}
\pagestyle{plain}

\begin{document}
{\Large\bfseries Closed-Loop and Activity-Guided Optogenetic Control\par}

\textbf{Paper:} Grosenick, Marshel and Deisseroth (2015).\quad Review of feedback, sensing and optical actuation.\par
\href{https://doi.org/10.1016/j.neuron.2015.03.034}{DOI: \nolinkurl{10.1016/j.neuron.2015.03.034}}\par
{\footnotesize Page references below use PDF page numbers, including front matter.}

\question{To which general question does this paper relate?}
What can be measured and controlled optogenetically, and how to formulate an experiment-specific control objective and model.

\question{How does it answer the general questions?}
\textbf{Framework.} Distinguishes feedback based on ongoing measurements from stimulation patterns chosen using earlier observations. Reviews electrical recordings, optical indicators and cell/projection-targeted light actuation, alongside on/off, PID, adaptive and predictive control approaches (PDF pp.~1--12).

\textbf{Usable insight.} The desired neural or behavioural trajectory determines sensor bandwidth, actuator choice and the meaning of successful control. Calcium signals do not necessarily resolve individual spike timing; opsin kinetics, saturation, targeting, heating and nonstationarity constrain achievable responses. Optogenetic actuation does not itself provide a neural measurement.

\textbf{Model relevance.} Supports system identification and explicit sensor/opsin dynamics before feedback design. A memoryless input map or time-invariant delayed rate model is a chosen approximation, not a property guaranteed by the technology. Examples include seizure interruption and manipulation of oscillations, but no universal beta-suppression objective or PIE Zames--Falb optimization is provided (PDF pp.~10--12, 24--26).

\question{What papers does it cite that could provide more information?}
\begin{itemize}
\item Paz et al. (2013), \emph{Closed-loop optogenetic control of thalamus as a tool for interrupting seizures after cortical injury} (references, PDF p.~32).
\item Krook-Magnuson et al. (2014), \emph{Cerebellar directed optogenetic intervention inhibits spontaneous hippocampal seizures in a mouse model of temporal lobe epilepsy} (PDF p.~30).
\end{itemize}

\end{document}
